\section{The optimum within the join type}
\label{sec:pair}

\begin{theorem}
  \label{thm:pair}
  Let $m, l \in \{3, 4, 5\}$, $p \in (\mathbb{S}^{3})^{m+l}$, and $D(p)$ the
  signed volume \eqref{eq:Djoin} of $\mathcal{J}_{m,l}$, with $a_{i}$, $b_{j}$
  as in Definition~\ref{def:join}. Then
  \begin{equation}
    \label{eq:pairbound}
    |D(p)| \leq m\sin\frac{2\pi}{m}\cdot l\sin\frac{2\pi}{l},
  \end{equation}
  with equality if and only if, for some $2$-plane $\Pi$, the points $a_{1},
  \ldots, a_{m}$ are the consecutive vertices of a regular $m$-gon inscribed in
  the unit circle of $\Pi$ and $b_{1}, \ldots, b_{l}$ those of a regular
  $l$-gon inscribed in the unit circle of $\Pi^{\perp}$.
\end{theorem}

\begin{proof}
  With $A$ and $B$ as in Definition~\ref{def:join},
  Lemma~\ref{lem:factor} gives $|D(p)| \leq |A|\,|B|$, and by
  Proposition~\ref{prop:polygon}, $|A| \leq m\sin\frac{2\pi}{m}$ and $|B| \leq
  l\sin\frac{2\pi}{l}$. This proves \eqref{eq:pairbound}.

  If equality holds, then $A, B \neq 0$ and equality holds in all three
  inequalities. By Proposition~\ref{prop:polygon}, the $a_{i}$ are as in
  \eqref{eq:regular} in a plane $\Pi$, with $A = m\sin\frac{2\pi}{m}\,
  \omega_{\Pi}$, and the $b_{j}$ in a plane $\Pi'$, with $B =
  l\sin\frac{2\pi}{l}\,\omega_{\Pi'}$. By Lemma~\ref{lem:factor}, $\star B =
  \lambda A$, so $\star\omega_{\Pi'}$ is a multiple of $\omega_{\Pi}$; both
  have length $1$, as $\star$ is an isometry, so by \eqref{eq:starplane}
  $\omega_{\Pi'^{\perp}} = \pm\star\omega_{\Pi'} = \pm\omega_{\Pi}$. With
  orthonormal bases $u_{1}, u_{2}$ of $\Pi$ and $w_{1},
  w_{2}$ of $\Pi'^{\perp}$, $\omega_{\Pi} = u_{1}u_{2}^{\mathsf{T}} -
  u_{2}u_{1}^{\mathsf{T}}$ has column space $\Pi$, and likewise for
  $\Pi'^{\perp}$; hence $\Pi = \Pi'^{\perp}$.

  Conversely, let $u_{1}, u_{2}$ and $u_{3}, u_{4}$ be orthonormal bases of
  $\Pi$ and $\Pi^{\perp}$ for which the $a_{i}$ and $b_{j}$ have the form
  \eqref{eq:regular}. By Proposition~\ref{prop:polygon}, $A = m\sin\frac{2\pi}{m}\,u_{1}
  \wedge u_{2}$ and $B = l\sin\frac{2\pi}{l}\,u_{3} \wedge u_{4}$, and by
  Lemma~\ref{lem:factor} and \eqref{eq:starplane}, $|D(p)| = |\langle A, \star
  B\rangle| = |A|\,|B|$.
\end{proof}

\begin{remark}
  \label{rem:bipyramid}
  The boundary complex of a double $n$-pyramid in $\mathbb{R}^{3}$ is the join
  of an $(n-2)$-cycle and a pair of points, and Berman and Hanes showed that an inscribed double
  $n$-pyramid with Property Z is unique up to congruence, with volume
  $\frac{n-2}{3}\sin\frac{2\pi}{n-2}$ \cite[Lemma~2]{berman-hanes-1970}; for
  $d$-bipyramids, Horv\'ath and L\'angi showed that Property Z forces the
  apexes to be antipodal and the base to lie in the orthogonal hyperplane
  \cite[Corollary~2]{horvath-langi-2016}. Theorem~\ref{thm:pair} is the
  analogue in $\mathbb{R}^{4}$ for $\mathcal{J}_{m,l}$, the join of two
  cycles. Unlike these
  results, which use Property Z, it bounds $|D(p)|$ for every
  $p \in (\mathbb{S}^{3})^{m+l}$.
\end{remark}

\begin{corollary}
  \label{cor:pair}
  Let $m, l \in \{3, 4, 5\}$. If $P \in \mathcal{P}_{4}(m+l)$ is
  combinatorially equivalent to the polytope of Lemma~\ref{lem:freesum}, then
  \[
    \vol_{4}(P) \leq \frac{1}{24}\, m\sin\frac{2\pi}{m}\; l\sin\frac{2\pi}{l}
      = \frac{1}{6}\area(R_{m})\area(R_{l}),
  \]
  with equality if and only if $P$ is the polytope of Lemma~\ref{lem:freesum}
  for some orthonormal basis $u_{1}, \ldots, u_{4}$, that is, up to an
  orthogonal transformation, the convex hull of $R_{m}$ and $R_{l}$ inscribed
  in $\mathbb{S}^{3}$ in orthogonal $2$-planes through the origin.
\end{corollary}

\begin{proof}
  Label the vertices of $P$ so that $p \in (\mathbb{S}^{3})^{m+l}$ realises
  $\mathcal{J}_{m,l}$. By Lemmas~\ref{lem:coherent} and~\ref{lem:volume},
  $24\vol_{4}(P) = |D(p)|$, and Theorem~\ref{thm:pair} gives the bound. In the
  case of equality, Theorem~\ref{thm:pair} and \eqref{eq:regular} give the
  basis $u_{1}, \ldots, u_{4}$. Conversely, by Lemma~\ref{lem:freesum} the
  polytope of that lemma lies in $\mathcal{P}_{4}(m+l)$ and its vertices, in
  the order given there, realise $\mathcal{J}_{m,l}$; they have the form of
  the equality case of Theorem~\ref{thm:pair}, so its volume is
  $\frac{1}{24}\, m\sin\frac{2\pi}{m}\; l\sin\frac{2\pi}{l}$.
\end{proof}
